\section{Relative linear nullhomotopies and low quadratic residues}
\label{sec:136-relative-disk-factorization}

\paragraph{Archive and exact scope.}
The uniform assumptions are cote136, physical pp.~117 and119,
footers116 and118, scans117 and119; the common-source question
is physical p.~121, footer120, scan121.  This module gives a
different reason for certain native diagonal residues to vanish.
It does not assume that the native quadratic homology is killed
by a coefficient tail.  That assertion is false by
ARCH-136-QUADRATIC-TAIL-ACTION.
The new conclusion requires the original native weight itself
to be in the target's bounded vertical comparison range.
Increasing only the coefficient-tail parameter does not establish
that hypothesis.

Throughout, \(R=\mathbb Z\) or \(\mathbb Z_{(p)}\),
\(S=\Gamma_R(T)\), and all linear maps preserve weight.
For a nonnegative cochain complex \(K\), write
\(\Gamma_{\mathrm{gr}}(K)\) for the free special strictly
commutative differential PD algebra on its graded module:
even generators have divided powers and odd generators have
square zero.  The differential is the special derivation
extending that of \(K\).  No freeness of \(K^0\) is
implicit in this notation.

\paragraph{ARCH-136-RELATIVE-DISK-FACTORIZATION.}
Suppose \(K\) is a nonnegative cochain complex, every
\(K^q\) for \(q>0\) is weight-graded free over \(S\),
and \(f:K\to C\) is a linear cochain map into a special
strict differential PD algebra.  Let \(J\subseteq C\)
be a differential PD ideal containing the image of \(f\).
Suppose there is a weight-preserving linear homotopy
\[
 h_q:K^q\longrightarrow J^{q-1}\quad(q\geq1),\qquad
 f_q=d h_q+h_{q+1}d_K,
 \quad f_0=h_1d_K.                                      \tag{136.RF1}
\]
Then the induced strict map from
\(\Gamma_{\mathrm{gr}}(K)\) factors through a free
nonnegative differential PD algebra \(E\) on unit disks,
and the augmentation ideal \(E_+\) maps into \(J\).
Moreover
\[
 H^0(E)=S,\qquad H^1(E)=H^2(E)=H^3(E)=0.                 \tag{136.RF2}
\]
These conclusions hold in each weight and separately in each
positive auxiliary arity.  Consequently, for every integer
\(N\geq1\), the image of every closed degree-two Moore-cone
pair in a positive-arity part of
\(\Gamma_{\mathrm{gr}}(K)\) has a cone primitive in
\(J\).  In particular that image has zero relative graph residue
when \(J=C(X,Y)\) is a restriction kernel.

\begin{proof}
For each homogeneous basis element \(x\) of \(K^q\),
\(q\geq1\), introduce \(u_q(x)\) in form degree
\(q-1\) and \(v_q(x)\) in degree \(q\), both of
the weight of \(x\), with
\[
                     d u_q(x)=v_q(x),\qquad d v_q(x)=0.
                                                               \tag{136.RF3}
\]
Extend \(u_q,v_q\) linearly on \(K^q\), and let
\(E\) be the free strict PD algebra on all these unit
disks.  Define a linear map \(j:K\to E_+\) by
\[
 \begin{aligned}
 j_0(x)&=u_1(d_Kx),\\
 j_q(x)&=v_q(x)+u_{q+1}(d_Kx)\qquad(q\geq1).
 \end{aligned}                                                \tag{136.RF4}
\]
It is defined even when \(K^0\) is nonfree, since only
its linear map to the free module \(K^1\) occurs.
Directly,
\(d j_q(x)=v_{q+1}(d_Kx)=j_{q+1}(d_Kx)\), with
the same identity for \(q=0\), because \(d_K^2=0\).
Thus \(j\) extends to a strict differential PD map.
There is a strict map \(E\to C\) taking
\[
                   u_q(x)\longmapsto h_q(x),\qquad
                   v_q(x)\longmapsto d h_q(x).
                                                               \tag{136.RF5}
\]
Its generators land in \(J\), and (136.RF1) says
exactly that its composite with \(j\) is \(f\).
The free universal properties prove the claimed factorization.

For completeness, the low-cohomology calculation does not use
flatness of \(K^0\), which is absent from the generating
modules of \(E\).  A disk whose lower generator has even
degree \(a\geq0\) has nonconstant blocks
\[
 S u^{[r]}\xrightarrow{1}S u^{[r-1]}v
                    \quad\text{in forms }ar,ar+1.
\]
Each is actually contractible.  A disk whose lower degree
\(a\geq1\) is odd has the nonconstant blocks
\[
 S u v^{[r-1]}\xrightarrow{r}S v^{[r]}
       \quad\text{in forms }(a+1)r-1,(a+1)r.
                                                               \tag{136.RF6}
\]
The \(r=1\) block is contractible.  For \(r\geq2\)
the first possible lower degree is three, and that lower
differential is injective over the integer-torsion-free ring
\(S\).  Its first possible cohomology is in degree four.
Decompose the tensor product of the disk algebras as a direct
sum of products of their constant and nonconstant blocks.
Any component containing a contractible block is contractible
by its actual signed tensor homotopy.  A component containing
at least two noncontractible blocks starts in degree at least
six; a single such block has no cohomology below degree four.
The same argument applies to infinitely many disks because
every tensor monomial has finite support.  It proves (136.RF2)
in each weight and auxiliary arity.  Assigning arity one to
all generators also makes (136.RF4) compatible with the source
arity on \(\Gamma_{\mathrm{gr}}(K)\).  The target need not
carry that auxiliary grading.

Use the cone convention
\(P_N(A)^r=A^r\oplus A^{r+1}\),
\(d(x,y)=(dx-Ny,-dy)\).  Its degree-two cohomology is an
extension of \(H^2(A)/N\) by \(H^3(A)[N]\).
This group is zero for each positive-arity part of \(E\).
A closed source cone pair therefore maps to a boundary in
that part of \(E\).  Map a primitive forward using
(136.RF5); it belongs to \(J\).  In a restriction kernel
its restriction to \(Y\) is literally zero, which proves
the final assertion.  No choice of a linear family of cone
primitives has been made or needed.
\end{proof}

\paragraph{The coefficient Ext input with a finite quotient.}
Put \(I=I_M\), \(M\geq1\), and
\(F_N=S/(N,I)\).  For every nonnegative weight shift
\(v\), ARCH-136-TAIL-EXT-VANISHING gives
\[
 \operatorname{Ext}^i_{S,\mathrm{gr}}
             ((S/I)\langle v\rangle,S)=0\quad(i\geq0).
\]
The quotient \(S/I\) is integer-torsion-free.  Applying
Hom to its multiplication-by-\(N\) sequence consequently gives
\[
 \operatorname{Ext}^i_{S,\mathrm{gr}}
              (F_N\langle v\rangle,S)=0\quad(i\geq0).
                                                               \tag{136.RF7}
\]
This includes the case that \(N\) is a unit, when the source
is zero.  The category is the nonnegative weight-graded category
used in that checked theorem; its free resolutions also compute
the same groups in all graded modules.

\paragraph{ARCH-136-HIGH-NATIVE-RESIDUE.}
Let \(D\geq4\) be even and \(v\geq0\).  Suppose a nonnegative complex
\(K\), concentrated in forms zero through \(D\), has
\[
 H^q(K)=0\ (q\ne D),\qquad
 H^D(K)=F_N\langle v\rangle,\qquad K^q\text{ free for }q>0,
                                                               \tag{136.RF8}
\]
and is a free resolution of that last module truncated at its
\(D\)-th syzygy, with every homogeneous generator having
weight at least \(v\).  Let \(C\) be admissible and
\(Y\subseteq X\) any simplicial pair.  There is a bound
\(W_C(D)\), independent of \(X,Y,M,N\) and of all ranks
in the resolution, such that
\[
                          v\geq W_C(D)                 \tag{136.RF9}
\]
implies that every linear map \(K\to C(X,Y)\) is
linearly nullhomotopic inside that restriction kernel.
It therefore satisfies ARCH-136-RELATIVE-DISK-FACTORIZATION.

For the actual native degree-four block of
ARCH-136-NATIVE-MIXED-TOR, one has
\(N=2\), \(D=4\), and \(v=2m\).  If
\(2m\geq W_C(4)\), every quadratic degree-three Moore
pair from that block has zero individual graph residue in
\(C\).  This includes the non-tail-annihilated diagonal
classes of ARCH-136-QUADRATIC-TAIL-ACTION.  A finite collection
of targets is covered by the maximum of their bounds.  The same
conclusion holds after a specified homogeneous RH1 stage whenever
its positive-form terms in that auxiliary arity are unchanged;
no arbitrary continuation is included in this statement.

\begin{proof}
Complete the reversed truncated resolution to a free homological
resolution \(P_\bullet\to F_N\langle v\rangle\).
For \(i<D\) put \(P_i=K^{D-i}\), and take a free cover
\(P_D\to K^0\), followed by free covers of its successive
kernels.  All generator weights remain at least \(v\),
since every module being covered has no smaller weight.
The quotient of \(P_D\) by the next image is precisely
\(K^0\).

The proof of ARCH-136-DERIVED-TAIL-COMPARISON applies verbatim
to this resolution.  Its comparison proof uses only freeness,
the common lower bound on generator weights, and the coefficient
complex \(\operatorname{Hom}_{S,\mathrm{gr}}(P_\bullet,S)\);
it never uses that the resolved module is the particular ideal
\(I_M\).  More explicitly, take \(W_C(D)\) to dominate
the vertical and horizontal bounds through form degree \(D+1\).
In total degrees at most \(D+1\), every Hom term is a product
of target rows in that range.  Products are exact, and the
nonnegative totalizations have finite diagonals.  The same
projective relative-section resolution and its recursive lifting
proof give
\[
 H^j\operatorname{Tot}\operatorname{Hom}_{S,\mathrm{gr}}
                  (P_\bullet,C(X,Y))
  \simeq H^j\bigl(X,Y;
                \operatorname{Hom}_{S,\mathrm{gr}}(P_\bullet,S)
             \bigr),\quad j\leq D.                    \tag{136.RF10}
\]
The coefficient complex on the right is acyclic by (136.RF7).
Relative normalized chains are a nonnegative complex of free
abelian groups, so their Hom totalization into that nonnegative
acyclic coefficient complex is acyclic: its first-quadrant
spectral sequence has zero vertical cohomology and finite
diagonals.  Thus the left side in degree \(D\) is zero.

A map \(f:K\to C(X,Y)\) gives a total degree-\(D\)
cocycle in the left complex.  In indices \(i<D\) use
\(f_{D-i}:P_i=K^{D-i}\to C^{D-i}(X,Y)\); in index
\(D\) use its composite with \(P_D\to K^0\), and
use zero above it.  With the convention
\[
       (\delta a)_i=d a_i-(-1)^{|a|}a_{i-1}\partial_i,
                                                               \tag{136.RF11}
\]
it is a cocycle because \(D\) is even and \(f\) is
a cochain map.  It is a boundary, say \(f=\delta h\).
The total degree of \(h\) is \(D-1\), so it has no
component in resolution degree \(D\) or higher.  Its
components in degrees \(0,\ldots,D-1\) yield
\(h_q:K^q\to C^{q-1}(X,Y)\) for \(q>0\), and
(136.RF11) gives exactly (136.RF1).  At the bottom,
the equation on \(P_D\) descends to
\(f_0=h_1d_K\), because \(P_D\to K^0\) is onto.
Every component already has values in the restriction kernel.
Apply the preceding factorization theorem.

For the native degree-four block, (136.RF8) and the generator
bound are exactly its checked truncated-resolution description.
In low forms the actual auxiliary-arity-four component is
\(\Gamma^2_{\mathrm{gr}}(K)\); the only native multiplication
factors differing from one occur in forms seven and eight,
outside the degree-two cone and its degree-one primitives.
Thus the established native quadratic comparison gives the
stated individual residue vanishing.  This argument neither
annihilates the source cohomology nor asserts that the source
quadratic class was already a boundary.  A stage leaving all
positive-form terms in the indicated arity unchanged leaves
the degree-two cone cycles and degree-one cone boundaries
unchanged, since these use only forms one through four.
\end{proof}

\paragraph{Exact consequence and remaining gate.}
For a fixed native form degree, all blocks whose original weight
lies above the stated target bound have zero nonlinear low-cone
residue through this factorization.  In particular the diagonal
degree-four case is controlled after imposing the one condition
\(2m\geq\max_C W_C(4)\).  This is distinct from increasing
the parameter-tail cutoff \(M\).  For native operations at
all primes, the analogous sufficient conditions are
\(pm\geq W_C(p(q+1))\); no finite upper bound for
\(W_C(p(q+1))/p\) is assumed or proved.  A horizontal
guard on the original input row also does not imply (136.RF9).
Low native weights, arbitrary subsequent old-algebra interactions,
and preservation of admissibility by further strict PD repairs
remain outside the result.  Integral connectedness remains open.

\paragraph{Verification status.}
The author has completed a full saved-proof reread, checking the
nonfree-bottom formula, every disk block and low cone degree,
the finite-quotient Ext sequence, the uniform relative-section
comparison, all homotopy signs and bottom descent, and the exact
native-weight limitation.  The parent and both peer branches have completed whole saved-proof
checks and reported PASS; their actual scopes and the root
confirmation are recorded separately.  No finite computation is used as a
substitute for the factorization or the Ext argument.
